% Part: second-order-logic
% Chapter: syntax-and-semantics
% Section: language-of-sol

\documentclass[../../../include/open-logic-section]{subfiles}

\begin{document}

\olfileid{sol}{syn}{lan}

\olsection{દ્વિતીય-ક્રમ તર્કશાસ્ત્રની ભાષા}

પ્રથમ-ક્રમ તર્કશાસ્ત્રની જેમ, દ્વિતીય-ક્રમ તર્કશાસ્ત્રના
વ્યંજકો મૂળભૂત શબ્દભંડોળમાંથી રચાય છે. તેમાં
\emph{!!{variable}s}, \emph{!!{constant}s},
\emph{!!{predicate}s} અને ક્યારેક \emph{!!{function}s}
હોય છે. તેમનાથી, તાર્કિક સંયોજકો, પરિમાણકો અને કૌંસ તથા
અલ્પવિરામ જેવાં વિરામચિહ્નો સાથે, \emph{પદો} અને
\emph{!!{formula}s} રચાય છે.
ફરક એટલો છે કે વસ્તુઓ માટેના ચલો ઉપરાંત દ્વિતીય-ક્રમ
તર્કશાસ્ત્રમાં સંબંધો અને ફલનો માટેના ચલો પણ છે,
અને તેમના પર પરિમાણીકરણની છૂટ છે.
એટલે દ્વિતીય-ક્રમ તર્કશાસ્ત્રનાં તાર્કિક ચિહ્નો
પ્રથમ-ક્રમનાં ચિહ્નો ઉપરાંત આ છે:

\begin{enumerate}
\item દરેક અરીટી~$n$ માટે દ્વિતીય-ક્રમ સંબંધ
!!{variable}sનો !!{denumerable}s ગણ:
$\Obj V_0^n$, $\Obj V_1^n$, $\Obj V_2^n$, \dots
\item દ્વિતીય-ક્રમ ફલન !!{variable}sનો !!{denumerable}s ગણ:
$\Obj u_0^n$, $\Obj u_1^n$, $\Obj u_2^n$, \dots
\end{enumerate}

પ્રથમ-ક્રમ તર્કશાસ્ત્રમાં સામાન્ય રીતે !!{identity}~$\eq$
સમાવાય છે. તેમાં ભાષા~$\Lang{L}$નાં અતાર્કિક ચિહ્નો
રસપ્રદ વાતો વ્યક્ત કરવા માટે અગત્યનાં છે.
અતાર્કિક ચિહ્નો વિનાનાં !!{sentences} પણ છે,
પરંતુ ફક્ત $\eq$થી રસપ્રદ વાતો કહેવી મુશ્કેલ છે.
દ્વિતીય-ક્રમ તર્કશાસ્ત્રમાં સંબંધ અને ફલનના ચલોનો અમર્યાદિત
પુરવઠો હોવાથી, વિશેષ અતાર્કિક ચિહ્નોને અનુરૂપ ચલો વડે
પ્રથમ-ક્રમ ભાષાની વાતો રજૂ કરી શકીએ;
અર્થ જાળવવા આ ચલોને તે ચિહ્નોનાં અર્થઘટનનાં મૂલ્યો આપવા પડે છે.
\sourcecorrection{OLMD-346}{મૂળના છેલ્લાં વાક્યમાં અતાર્કિક ચિહ્નો વિના પ્રથમ-ક્રમ ભાષાની કોઈ પણ વાત કહી શકાય તેવો અમર્યાદિત દાવો છે. સંબંધ અને ફલનના મુક્ત ચલો તે ચિહ્નોની જગ્યાએ મૂકી, તેમના અર્થઘટનને અનુરૂપ નિયુક્તિ આપતાં સંતોષ જાળવાય છે; અચળ માટે વસ્તુ-ચલ વાપરી શકાય છે. આ ચલો પર આપમેળે પરિમાણીકરણ કરી દેવાથી મૂળનું ચોક્કસ અર્થઘટન જળવાતું નથી. તેથી મુક્ત ચલો અને નિયુક્તિની જરૂરી શરત સ્પષ્ટ કરી છે.}

\end{document}

